\documentclass[11pt]{article}
\usepackage[a4paper,margin=1in]{geometry}
\usepackage{amsmath,amssymb,amsthm,mathtools}
\usepackage{enumitem}
\usepackage[colorlinks=true,linkcolor=blue,citecolor=blue,urlcolor=blue]{hyperref}
\usepackage{xurl}

\newcommand{\indep}{\mathrel{\perp\mkern-9mu\perp}}
\newcommand{\Renc}{\mathcal{R}_{\mathrm{enc}}}
\newcommand{\epscpa}{\varepsilon_{\mathrm{cpa}}}
\newcommand{\HunpC}{H_{\mathrm{unp}}^{\mathcal{C}}}
\newcommand{\Hunp}{H_{\mathrm{unp}}}
\newcommand{\C}{\mathcal{C}}
\newcommand{\Yspace}{\mathcal{Y}}
\newcommand{\Wspace}{\mathcal{W}}
\newcommand{\Encf}{\mathrm{Enc}}
\newcommand{\Decf}{\mathrm{Dec}}
\newcommand{\KeyGen}{\mathrm{KeyGen}}
\newcommand{\AliceView}{\mathit{AliceView}}
\newcommand{\AliceViewHyb}{\mathit{AliceView}_{\textrm{after-swaps}}'}
\newcommand{\Dk}{\mathit{Dk}}
\newcommand{\pk}{\mathit{pk}}

\title{$N$-party DSDP per-target secrecy}
\author{}
\date{}

\begin{document}
\maketitle

\section*{Setup}

This note uses the notation and computational assumptions from the 3-party Alice note, generalized from the key registry $(\pk_a, \pk_b, \pk_c)$ to an $N$-party public-key registry with one decryption key per party. The adversary class $\C$, real-or-zero indistinguishability parameter $\epscpa$, conditional unpredictability entropy $\HunpC$, and fresh encryption-randomness convention are unchanged.

\section*{Generalization to N parties}

The 3-party Alice theorem from the companion note generalizes to the original $N$-party DSDP protocol with one corrupted party once the residual-view shape is verified for the corrupted role. The proof structure is the same hybrid-then-fiber pattern, with a per-target leakage budget that scales linearly in $N$.

\subsection*{Per-target bound}

Let $N \geq 3$ be the number of parties, and suppose party $p \in \{1, \ldots, N\}$ is corrupted. The corrupted party's view contains
\begin{itemize}
\item the public protocol material such as opened sums and masks,
\item the corrupted party's own decryption key $\Dk_p$,
\item one ciphertext encrypted under each of the other $N-1$ parties' public keys.
\end{itemize}
Repeating the path-(b) hybrid argument of the companion 3-party theorem with $N-1$ cross-key ciphertext swaps in place of the two used in the 3-party Alice case gives the bound. Suppose additionally that there exists $j_0 \in \{1, \ldots, N\} \setminus \{p, i\}$ with $U_{j_0} \in R^\times$ almost surely. Then for every $A \in \C$ and every target $V_i$ with $i \neq p$,
\[
\Pr[A(\mathit{View}_p) = V_i] \leq \frac{1}{m} + (N-1) \cdot \epscpa.
\]
The hypothesis $U_{j_0} \in R^\times$ lets the residual constraint $\sum_{j \notin \{p, i\}} U_j V_j = c - U_i V_i$ be solved uniquely for $V_{j_0}$ given any candidate value of $V_i$. Each such candidate is then paired with exactly $m^{N-3}$ joint values of the remaining hidden variables, giving a fiber whose marginal in $V_i$ is uniform on $R$. The witness-unit condition is sufficient but not strictly necessary. Without such a witness coordinate, $V_i$ may be a determined coordinate of the fiber rather than a free one, in which case the per-target uniformity claim may fail. Whether it fails or holds in a particular setting depends on the ring and on the joint distribution of the coefficients.

The $(N-1) \cdot \epscpa$ term is the cost of $N-1$ applications of real-or-zero indistinguishability in the hybrid sequence $H_0 \to H_1 \to \cdots \to H_{N-1}$, one $\epscpa$ per swap.

\subsection*{Concrete security as $N$ grows}

The bound is linear in $N$. The standard concrete-security framing keeps it meaningful in the regime where the encryption advantage $\epscpa$ is negligible in a security parameter $\lambda$ and the number of parties is polynomially bounded in $\lambda$:
\[
\epscpa = \mathrm{negl}(\lambda), \qquad N = \mathrm{poly}(\lambda).
\]
Under these conditions,
\[
(N-1) \cdot \epscpa = \mathrm{poly}(\lambda) \cdot \mathrm{negl}(\lambda) = \mathrm{negl}(\lambda),
\]
so the per-target bound remains negligibly close to the IT baseline $1/m$ for every polynomially bounded $N$.

A concrete numerical example clarifies the regime. Take $\epscpa = 2^{-128}$, the standard 128-bit security level, and $m = 2^{32}$. The IT baseline is $1/m = 2^{-32}$. For $N = 10^6$ parties,
\[
(N-1) \cdot \epscpa \approx 2^{20} \cdot 2^{-128} = 2^{-108},
\]
which is roughly $2^{76}$ times smaller than $1/m$. The cryptographic term is dominated by the IT baseline, and the bound is essentially $1/m$. The crossover point at which the cryptographic term begins to compete with the IT baseline is $N \approx 1/(m \cdot \epscpa) = 2^{96}$, far beyond any realistic instantiation.

\subsection*{Per-target uncertainty does not grow with $N$}

The per-target uncertainty of $V_i$, measured by $\log m$ bits, is fixed in $N$. Adding parties to the protocol does not increase the entropy of any individual secret. The per-target bound therefore compares a fixed IT baseline $1/m$ against a leakage budget $(N-1) \cdot \epscpa$ that grows linearly in $N$. The protocol remains secure as long as the leakage stays below the IT baseline, that is,
\[
(N-1) \cdot \epscpa < \frac{1}{m}.
\]
For $\epscpa = \mathrm{negl}(\lambda)$ and $m = \mathrm{poly}(\lambda)$, this inequality holds whenever $N$ is super-polynomially smaller than $1/(m \cdot \epscpa) = 1/\mathrm{negl}(\lambda)$, in particular whenever $N$ is polynomial in $\lambda$.

\subsection*{Joint security across all targets}

A joint claim about all $N-1$ hidden secrets is bounded relative to the constraint imposed by the publicly opened sum $S$. Conditioning on $S$ and on the corrupted party's input $V_p$ reduces the IT support of $(V_j)_{j \neq p}$ from $m^{N-1}$ to $m^{N-2}$ elements: it is a coset of the kernel of the linear form $\sum_{j \neq p} U_j V_j = S - U_p V_p$, well-defined as soon as at least one $U_j$ with $j \neq p$ is a unit. The natural joint baseline is therefore $1/m^{N-2}$, and the corresponding hybrid bound is
\[
\Pr\big[ A \text{ guesses all of } (V_j)_{j \neq p} \text{ correctly} \big] \leq \frac{1}{m^{N-2}} + (N-1) \cdot \epscpa.
\]
For the joint bound to remain meaningful, $\epscpa$ must be smaller than $1/m^{N-2}$, that is, exponentially small in $(N-2) \cdot \log m$. For small $N$ the requirement is mild: at $N = 3$ it reduces to $\epscpa < 1/m$, which standard 128-bit IND-CPA schemes already satisfy when $m$ is polynomial in $\lambda$. As $N$ grows, the requirement becomes exponentially small in $N \log m$ and is unreachable with standard cryptographic schemes for large $N$, so joint security across all secrets is usually not claimed in that regime. Per-target security, with the linear-in-$N$ bound above, is the natural notion to maintain across all $N$.

\subsection*{Asymptotic per-target theorem statement}

The following is the precise asymptotic statement that the path-(b) reformulation of the formal development is intended to carry once the 3-party case is discharged. The current Coq formalization in \texttt{dsdp\_security.v} proves a Shannon-entropy statement under the unsound encryption-independence axiom. The asymptotic $N$-party Hunp/RoR theorem below is the target of that reformulation, not a proven result of the existing development.

\textbf{Theorem (Asymptotic per-target security under $N$-party scaling).} Fix a security parameter $\lambda$, a corrupted party $p \in \{1, \ldots, N\}$, and a target index $i \neq p$. Assume the following.
\begin{itemize}
\item \emph{Encryption.} The scheme satisfies real-or-zero indistinguishability against $\C$ with advantage $\epscpa = \epscpa(\lambda)$ negligible in $\lambda$. The size bound for $\C$ is large enough to absorb a hybrid-step wrapper $t_{\mathrm{hyb}}(N)$ of size polynomial in $N$ around any predictor used in the proof, yielding the per-step distinguishers. Since $N = N(\lambda)$ is polynomial, the $N$-step hybrid produces $N - 1$ distinguishers each of size at most $s_0 + t_{\mathrm{hyb}}(N) = \mathrm{poly}(\lambda)$, all within $\C$.
\item \emph{Scale parameters.} The number of parties is $N = N(\lambda)$ polynomially bounded in $\lambda$, and the message-space size is $m = m(\lambda)$ polynomially bounded in $\lambda$.
\item \emph{Input distribution.} All party inputs $(V_j)_{j=1, \ldots, N}$, including the corrupted party's input $V_p$, are jointly independent and individually uniform on $R$. In particular, the hidden secrets $(V_j)_{j \neq p}$ are jointly independent of $V_p$ and of any corrupted-party initial state or public-setup variable that appears in $\mathit{View}_p$.
\item \emph{Coefficient distribution.} There exists a witness index $j_0 \in \{1, \ldots, N\} \setminus \{p, i\}$. The witness coefficient $U_{j_0}$ is supported almost surely on $R^\times$, with otherwise arbitrary distribution. The remaining coefficients $(U_j)_{j \neq j_0}$ are individually uniform on $R$. The protocol-internal randomness $(R_j)$ is individually uniform on $R$. All coefficients $(U_j)_{j=1, \ldots, N}$ and randomness $(R_j)$ are jointly independent of $(V_j)_{j=1, \ldots, N}$ and of each other. In particular, $U_{j_0}$ is jointly independent of $(U_j)_{j \neq j_0}$ and of $(R_j)$, despite its non-uniform support.
\item \emph{Residual view shape (path-(b) contraction).} After replacing every cross-key ciphertext in $\mathit{View}_p$ by an encryption of $0_R$ and rewriting every own-key ciphertext as its decrypted plaintext, the resulting residual view is a deterministic function of: the public protocol setup material, the corrupted-party data $(V_p, \Dk_p)$, the coefficients $(U_j)_{j=1, \ldots, N}$ and protocol-internal randomness $(R_j)$, the publicly opened sum $S = \sum_{j=1}^{N} U_j \cdot V_j$, and fresh encryption randomness. No function of the hidden secrets $(V_j)_{j \neq p}$, linear or nonlinear, is revealed in the residual view except through $S$. In particular, no individual product $U_j V_j$, partial sum, or per-pair quantity involving $V_j$ for $j \neq p$ appears in the residual view.
\item \emph{Encryption randomness.} The encryption randomness for each ciphertext in $\mathit{View}_p$ is fresh, individually uniform on $\Renc$, and jointly independent of all input, coefficient, and protocol-internal randomness variables.
\end{itemize}
Then for every $A \in \C$,
\[
\Pr[A(\mathit{View}_p) = V_i] \leq \frac{1}{m} + (N-1) \cdot \epscpa,
\]
and the right-hand side equals $1/m + \mathrm{negl}(\lambda)$.

The corresponding entropy form, obtained by applying $-\log$ to the per-target guessing-probability bound and using the definition of $\HunpC$, is
\[
\HunpC(V_i \mid \mathit{View}_p)
\geq -\log\left( \frac{1}{m} + (N-1) \cdot \epscpa \right)
= \log m - \log\left(1 + (N-1) \, m \, \epscpa\right).
\]
For $\epscpa = \mathrm{negl}(\lambda)$, $N = \mathrm{poly}(\lambda)$, and $m = \mathrm{poly}(\lambda)$, the second term $\log(1 + (N-1) \, m \, \epscpa)$ is $\mathrm{negl}(\lambda)$, so
\[
\HunpC(V_i \mid \mathit{View}_p) \geq \log m - \mathrm{negl}(\lambda).
\]

\paragraph{Remaining proof obligation.} The Residual view shape hypothesis packages the bulk of the protocol-specific work that this asymptotic theorem leaves open. Discharging it for a concrete corrupted role $p$ amounts to verifying that the path-(b) contraction of $\mathit{View}_p$ has the stated functional form. The companion 3-party Alice note verifies this for the 3-party case with $p = a$ (Alice) by exhibiting the explicit residual $\AliceViewHyb$. The Bob and Charlie roles in the 3-party case, and arbitrary corrupted roles in the $N$-party case, require their own role-specific verification. The asymptotic theorem above is therefore a target statement: it is sound under the hypotheses listed, and it reduces the remaining formal work to producing the corresponding explicit residual view for each role.

\section*{Reference Notes}

Reference \cite[Definition~7]{hsiao-lu-reyzin} introduces conditional unpredictability entropy with the adversary class fixed concretely as polynomial-size circuits. This $N$-party note uses that viewpoint for a target theorem whose adversary class must absorb polynomial-in-$N$ hybrid wrappers. Reference \cite[Secs.~1.3.3 and 3.2.4]{goldreich-foc} is a textbook source for the non-uniform circuit model and for treating polynomial overheads inside circuit classes.

References \cite{bellare-concrete,shoup-games,blanchet-pointcheval-games} justify the game-hopping proof style used in the $N$-party statement: the proof pays one real-or-zero indistinguishability loss for each cross-key ciphertext swap. The protocol-specific work not discharged by this abstract theorem is the residual-view contraction for each corrupted role; once that contraction is proved, the residual information-theoretic step is the linear-fiber argument described in the theorem statement.

\begin{thebibliography}{9}\bibitem{hsiao-lu-reyzin}
Hsiao, Lu, Reyzin.
\emph{Conditional Computational Entropy, or Toward Separating Pseudoentropy from Compressibility}.
Advances in Cryptology -- Eurocrypt 2007, Lecture Notes in Computer Science 4515, pages 169--186.
\url{https://www.iacr.org/archive/eurocrypt2007/45150169/45150169.pdf}

\bibitem{bellare-concrete}
Bellare, Desai, Jokipii, Rogaway.
\emph{A Concrete Security Treatment of Symmetric Encryption}.
Full version of FOCS 1997 paper, September 2000.
\url{https://cseweb.ucsd.edu/~mihir/papers/sym-enc.pdf}

\bibitem{shoup-games}
Shoup.
\emph{Sequences of Games: A Tool for Taming Complexity in Security Proofs}.
January 18, 2006.
\url{https://www.shoup.net/papers/games.pdf}

\bibitem{blanchet-pointcheval-games}
Blanchet, Pointcheval.
\emph{Automated Security Proofs with Sequences of Games}.
Advances in Cryptology -- Crypto 2006, Lecture Notes in Computer Science 4117, pages 537--554.
\url{https://bblanche.gitlabpages.inria.fr/publications/BlanchetPointchevalCrypto06.pdf}

\bibitem{goldreich-foc}
Goldreich.
\emph{Foundations of Cryptography, Volume I: Basic Tools}.
Cambridge University Press, 2001.
\url{https://www.wisdom.weizmann.ac.il/~oded/foc-vol1.html}

\bibitem{goldreich-uniform}
Goldreich.
\emph{A Uniform-Complexity Treatment of Encryption and Zero-Knowledge}.
Journal of Cryptology, Vol.~6, No.~1, pages 21--53, 1993; revised version, July 1998.
\url{https://www.wisdom.weizmann.ac.il/~oded/PSX/uniform.pdf}

\end{thebibliography}

\end{document}
